\scanpage{42}
\mspage{129}

\deleted{et comme \unsure{tantôt}, on a
\[
p^*R^1f_*(F)\xrightarrow{\sim}R^1 f'_*(q^*F),
\]
puisque les deux $R^1$ se calculent fibre par fibre.
On trouve donc
\[
H^1(X,F)\simeq H^0(Y',R^1f'_*(q^*F))
\]
\unsure{Or}}

On a \unsure{un hom|une ligne|une liaison} de suites exactes d'ens
ponctués
\diagram{\scriptsize\begin{tikzcd}[ampersand replacement=\&,row sep=small,column sep=small,execute at end picture={\draw ([xshift=-1pt]lasttop.north west) -- ([xshift=-5pt]lasttop.north west) -- ([xshift=-5pt]lastbot.south west) -- ([xshift=-1pt]lastbot.south west);\draw ([xshift=1pt]lasttop.north east) -- ([xshift=5pt]lasttop.north east) -- ([xshift=5pt]lastbot.south east) -- ([xshift=1pt]lastbot.south east);}]
\deleted{\unsure{\leftarrow}}\,H^1(Y,f_*(F))\arrow[r,"u"]\arrow[d,"\alpha"]
\&H^1(X,F)\arrow[r,"v"]\arrow[d,"\beta"]
\&\overset{\unsure{\ni|\in}\,\eta}{H^0(Y,R^1 f_*(F))}\arrow[r,"w"]\arrow[d,"\gamma"]
\&|[alias=lasttop]|H^2(Y,f_*F)\arrow[d,"\delta"']\\
H^1(Y',f'_*(F'))\arrow[r,"u'"]
\&\underset{\unsure{\ni|\downarrow}\,\unsure{\xi'|\zeta'}}{H^1(X',F')}\arrow[r,"v'"]
\&\underset{\downarrow\,\eta'}{H^0(Y',R^1 f'_*(F'))}\arrow[r,"w'"]
\&|[alias=lastbot]|H^2(Y',f'_*F')
\end{tikzcd}}
% slanted
\margin{\unsure{où on a posé} $F'=q^*(F)$}

où les faisceaux $f_*F$, $R^1 f_*(F)$ sont loc constants
par hyp. et se calculent fibre par fibre, les
faisceaux $f'_*(F')$ et $R^1 f'_{\unsure{*|\circ}}(F')$ aussi, ce qui implique
que ces derniers sont can. iso. (par \unsure{les|ces} flèches de
chgt. de base) aux images inverses de leurs
homologues $f_*(F)$, $R^1 f_*(F)$. Comme $p$ est
une Hot-équiv, cela implique que les flèches
$\alpha,\gamma$ sont iso. Donc si $\xi'\in H^1(X',F')$, son image $\eta'$
dans $H^0(Y',R^1 f'_*(F'))$ provient d'un $\eta\in H^0(Y,R^1 f_*(F))$
bien déterminé. Il faut voir que cet $\eta$ provient
d'un $\xi\in H^1(X,F)$. Ceci (si je me rappelle bien)
est possible ssi une certaine obstruction $w(\eta)$
dans $H^2(Y,f_*F)$ est dans le sous-ensemble neutre des
$H^2$ (\uline{NB} \deleted{\unsure{La}} Le champ \inserted{sur $\widehat Y$} des $F$-torseurs sur \inserted{les} $\widehat X/U$
qui remontent $\eta$ est une \deleted{\unsure{ob}} $f_*(F)$ gerbe, \inserted{$\mathscr G_\eta$} et il faut
voir qu'elle est neutre. Mais il y a mieux à faire :

% archive sheet number: 42
